\clearpage
\section{Related work}
Schuld and Killoran connected quantum learning models with feature maps into Hilbert spaces, formalizing quantum circuits as kernel-style data transformations followed by a measurement-based linear hypothesis [1]. Pérez-Salinas et al. demonstrated that repeatedly encoding classical inputs through multiple circuit layers---termed \emph{data re-uploading}---increases the expressivity of small quantum classifiers, establishing that even single-qubit systems with sufficiently many re-uploading stages can approximate arbitrary continuous functions [2].

Subsequent mathematical foundations revealed that variational quantum circuits act effectively as truncated multidimensional Fourier series [4]:
\begin{equation}
  f(x;\theta) = \sum_{\omega \in \Omega} c_\omega(\theta) e^{i \omega^\top x},
\end{equation}
where the frequency spectrum $\Omega$ is determined by the spectrum of the encoding generators and the number of re-uploading repetitions. For an $R_y(x)$ encoding layer, each repetition adds integer harmonic frequencies $\{-1, 0, 1\}$. Thus, shallow QNNs operate under strict spectral frequency cutoffs, unlike classical MLPs whose piecewise-linear ReLU activations can compose high-frequency boundary segments with localized support.

Bowles, Ahmed, and Schuld evaluated a wide range of variational quantum classifiers across tabular datasets, demonstrating that standard classical baselines (such as ridge classifiers and shallow neural networks) systematically outperform QML models [3]. Crucially, they observed that quantum entanglement frequently yielded negligible performance improvements and in several instances degraded generalization. Furthermore, McClean et al. proved that random parameterized quantum circuits suffer from the \emph{barren plateau} phenomenon [5], wherein gradient variances vanish exponentially with qubit count ($\text{Var}[\partial_\theta \mathcal{L}] \in \mathcal{O}(2^{-n})$). While our two-qubit investigation avoids exponential barren plateaus, it exposes the subtle interplay between circuit depth, entanglement, and inductive bias.

To address the historical scarcity of standardized quantum machine learning benchmarks, Perrier, Youssry, and Ferrie introduced \emph{QDataSet} [6], a comprehensive open-access repository of quantum datasets specifically curated for quantum machine learning research. QDataSet provides experimentally acquired and simulated quantum dynamics, state tomography data, and quantum state classification tasks across multi-qubit systems. As discussed in Section~\ref{sec:frontiers}, such quantum-native datasets constitute the primary regime where quantum neural networks possess a fundamental advantage over classical architectures by bypassing classical measurement bottlenecks.

\begin{table}[H]
\centering\small
\caption{How prior foundational literature informs our experimental and theoretical framework.}
\begin{tabularx}{\linewidth}{@{}p{0.24\linewidth}p{0.34\linewidth}X@{}}
\toprule
Prior direction & Core principle & Role in this study\\
\midrule
Quantum feature maps [1] & Unitary feature encoding maps classical inputs into Hilbert space inner products. & Two normalized features mapped to rotation angles on a 2-qubit register.\\
Data re-uploading [2] & Interleaving data encoding and trainable layers expands reachable hypothesis classes. & Compares depths 1, 2, and 3 to measure expressivity gains versus parameter costs.\\
Fourier representation [4] & VQCs compute trigonometric Fourier polynomials with bounded frequency spectra. & Explains QNN rotational advantages on circles and boundary limitations on Moons.\\
Benchmarking rigor [3] & Classical baselines routinely dominate small tabular tasks; entanglement benefits are task-contingent. & Uses matched splits, wide and compact MLPs, paired tests, and entanglement ablations.\\
Barren plateaus [5] & Expressivity-trainability trade-offs severely limit overparameterization in quantum circuits. & Theoretical basis for why classical MLPs benefit from overparameterization while QNNs cannot.\\
QDataSet benchmarks [6] & Standardized quantum state and dynamics datasets for benchmarking quantum algorithms. & Establishes quantum-native data standards where QNNs eliminate tomography overheads.\\
\bottomrule
\end{tabularx}
\end{table}
